% Options for packages loaded elsewhere
\PassOptionsToPackage{unicode}{hyperref}
\PassOptionsToPackage{hyphens}{url}
%
\documentclass[
  twocolumn]{article}
\usepackage{lmodern}
\usepackage{amssymb,amsmath}
\usepackage{ifxetex,ifluatex}
\ifnum 0\ifxetex 1\fi\ifluatex 1\fi=0 % if pdftex
  \usepackage[T1]{fontenc}
  \usepackage[utf8]{inputenc}
  \usepackage{textcomp} % provide euro and other symbols
\else % if luatex or xetex
  \usepackage{unicode-math}
  \defaultfontfeatures{Scale=MatchLowercase}
  \defaultfontfeatures[\rmfamily]{Ligatures=TeX,Scale=1}
\fi
% Use upquote if available, for straight quotes in verbatim environments
\IfFileExists{upquote.sty}{\usepackage{upquote}}{}
\IfFileExists{microtype.sty}{% use microtype if available
  \usepackage[]{microtype}
  \UseMicrotypeSet[protrusion]{basicmath} % disable protrusion for tt fonts
}{}
\makeatletter
\@ifundefined{KOMAClassName}{% if non-KOMA class
  \IfFileExists{parskip.sty}{%
    \usepackage{parskip}
  }{% else
    \setlength{\parindent}{0pt}
    \setlength{\parskip}{6pt plus 2pt minus 1pt}}
}{% if KOMA class
  \KOMAoptions{parskip=half}}
\makeatother
\usepackage{xcolor}
\IfFileExists{xurl.sty}{\usepackage{xurl}}{} % add URL line breaks if available
\IfFileExists{bookmark.sty}{\usepackage{bookmark}}{\usepackage{hyperref}}
\hypersetup{
  hidelinks,
  pdfcreator={LaTeX via pandoc}}
\urlstyle{same} % disable monospaced font for URLs
\usepackage[left=0.5cm,right=0.5cm,top=1.5cm,bottom=1.2cm]{geometry}
\usepackage{graphicx}
\makeatletter
\def\maxwidth{\ifdim\Gin@nat@width>\linewidth\linewidth\else\Gin@nat@width\fi}
\def\maxheight{\ifdim\Gin@nat@height>\textheight\textheight\else\Gin@nat@height\fi}
\makeatother
% Scale images if necessary, so that they will not overflow the page
% margins by default, and it is still possible to overwrite the defaults
% using explicit options in \includegraphics[width, height, ...]{}
\setkeys{Gin}{width=\maxwidth,height=\maxheight,keepaspectratio}
% Set default figure placement to htbp
\makeatletter
\def\fps@figure{htbp}
\makeatother
\setlength{\emergencystretch}{3em} % prevent overfull lines
\providecommand{\tightlist}{%
  \setlength{\itemsep}{0pt}\setlength{\parskip}{0pt}}
\setcounter{secnumdepth}{-\maxdimen} % remove section numbering
\usepackage{fancyhdr} \usepackage{booktabs,xcolor} \pagestyle{fancy} \renewcommand{\headrulewidth}{0pt} \rhead{2020} \cfoot{\includegraphics[width=20cm]{footer.png}} \fancypagestyle{plain}{\pagestyle{fancy}} \setlength{\headheight}{77.3pt} \setlength{\footskip}{36.5pt} \setlength{\textheight}{0.90\textheight} \pagenumbering{gobble} \usepackage[fontsize=9pt]{scrextend} \usepackage{float} \restylefloat{table} \usepackage{xcolor} \usepackage{multicol} \usepackage{array} \usepackage{colortbl} \usepackage{multirow} \usepackage{collcell} \usepackage{setspace} \usepackage{arydshln} \usepackage{caption} \captionsetup{skip=0pt} \setlength{\columnsep}{1.5cm}
\ifluatex
  \usepackage{selnolig}  % disable illegal ligatures
\fi

\author{}
\date{\vspace{-2.5em}}

\begin{document}

\newcommand{\greynote}[1]{
    {\scriptsize
    \textcolor{darkgray}{\textit{Notes:} #1}
  }
}

\newcommand{\greysource}[1]{
    {\scriptsize
    \textcolor{darkgray}{\textit{Source:} #1}
  }
}

\newcommand{\greydisclaimer}[1]{
    {\scriptsize
    \textcolor{darkgray}{\textit{Disclaimer:} #1}
  }
}

\newcommand*{\tabindent}{\hspace{1mm}}

\hypertarget{introduction}{%
\subsubsection{\texorpdfstring{\textbf{INTRODUCTION}}{INTRODUCTION}}\label{introduction}}

\textbf{The Global Education Policy Dashboard (GEPD): An innovative tool
to measure drivers of learning outcomes in basic education}\\
GEPD uses 3 data collection instruments to report on nearly 40
indicators that operationalize the World Development Report 2018
framework to track 3 areas for progress in education- Practices,
Policies, and Politics. Using these indicators, the dashboard highlights
areas where countries need to act to improve learning outcomes and
allows a way for governments to track progress as they act to close gaps
in these areas. For more information on GEPD, please visit
\textbf{\href{https://www.worldbank.org/en/topic/education/brief/global-education-policy-dashboard}{www.worldbank.org/global-education-policy-dashboard}}

\hypertarget{figure-1.-gepd-framework-practices-policies-and-politics-expanding-on-wdr-2018-framework}{%
\paragraph{Figure 1. GEPD framework (practices, policies and politics),
expanding on WDR 2018
framework}\label{figure-1.-gepd-framework-practices-policies-and-politics-expanding-on-wdr-2018-framework}}

\begin{center}\includegraphics[width=0.8\linewidth]{C:/Users/wb469649/Documents/Github/GEPD/Country_Reports/Data/full_circle} \end{center}

\hypertarget{instruments-of-data}{%
\subsubsection{\texorpdfstring{\textbf{INSTRUMENTS OF
DATA}}{INSTRUMENTS OF DATA}}\label{instruments-of-data}}

The \textbf{School Survey} consists of 8 modules to collect data across
200-300 schools on practices (the quality of service delivery in
schools) and de facto policy indicators. It consists of streamlined
versions of existing instruments together with new questions to fill
gaps in those instruments.\\
The \textbf{Policy Survey} collects information via interviews with
\textasciitilde200 officials per country at federal and regional level
to feed into the policy de jure indicators and identify key elements of
the policy framework.\\
The \textbf{Survey of public officials} collects information about the
capacity and orientation of the bureaucracy and political factors
affecting education outcomes. This survey is an education-focused
version of the civil-servant surveys from the Bureaucracy Lab, WBG.

\hypertarget{key-takeaways-2020}{%
\subsubsection{\texorpdfstring{\textbf{KEY TAKEAWAYS,
\uppercase{Jordan},
2020}}{KEY TAKEAWAYS, , 2020}}\label{key-takeaways-2020}}

\begin{itemize}
\tightlist
\item
  52\% learning poverty is observed in Grade 4. GEPD Grade 4 student
  proficiency (\textgreater80\% student knowledge) is only 4\%,
  attributed to low numeracy proficiency 2\%.
\item
  Low teacher content knowledge attributed to poor teacher support and
  lack of strong instructional leadership (especially principal feedback
  on classroom observation)
\item
  Low capacity for learning in Grade 1 (particularly socio-emotional
  knowledge) is attributed to low enrolment in early childhood programs
  (13\%)
\item
  A \textasciitilde0.8 point difference exists in the policy de-jure and
  de-facto indicators on a 5 point scale for teaching and school
  management. Major gaps are seen in teacher support, teacher motivation
  and evaluation of school management.
\item
  Quality of teaching support emerges as a major area requiring
  intervention, particularly through in-classroom support/feedback and
  pre-service training.
\item
  Bureaucracy scores high on mandates and accountability but low on
  impartial decision making
\end{itemize}

\vfill

\setlength\dashlinedash{0.2pt}
\setlength\dashlinegap{1.5pt}
\setlength\arrayrulewidth{0.3pt}

\hypertarget{table-1.-key-gepd-outcome-indicators}{%
\paragraph{Table 1. Key GEPD Outcome
Indicators}\label{table-1.-key-gepd-outcome-indicators}}

\begin{table}[H]
\resizebox{\columnwidth}{!}{\begin{tabular}{m{8cm}c}
\textbf{Indicator} & \textbf{Value}  \\\hline
Learning poverty & {\cellcolor{red!15}52\%} \\\cdashline{1-2}
Proficiency by End of Primary & {\cellcolor{red!15}50\%} \\\cdashline{1-2}
Proficiency on GEPD Assessment & {\cellcolor{red!15}4.4\%} \\\cdashline{1-2}
\hspace{1mm}\emph{Literacy proficiency} & {\cellcolor{red!15}22\%} \\\cdashline{1-2}
\hspace{1mm}\emph{Numeracy proficiency} & {\cellcolor{red!15}2\%} \\\cdashline{1-2}
Proficiency by Grade 2/3 & {NA} \\\cdashline{1-2}
Net Adjusted Enrollment Rate & {NA} \\\hline
\end{tabular}}
\\
\setstretch{0.8}\color{darkgray}\scriptsize{\textit{Source:} UIS, GLAD, GEPD, World Bank, Jordan, 2020. For information on indicators, please consult the World Bank \href{https://github.com/worldbank/GEPD}{\underline{GEPD}}, \href{https://github.com/worldbank/GLAD}{\underline{GLAD}} and \href{https://github.com/worldbank/LearningPoverty}{\underline{Learning Poverty}} repositories.}\\
\setstretch{0.8}\color{darkgray}\scriptsize{\textit{Notes:} (1) Proficiency on GEPD assessment means \% students with knowledge\textgreater{80\%}. (2) Proficiency by end of primary uses threshold as per Minimum Proficiency Levels set by GAML(UIS). (3) All indicators are on a scale of 0-5 unless measured in \%. (4) Green indicates indicator 'on-target', yellow indicates 'requires caution', red indicates 'needs improvement'.}
\end{table}
\raggedbottom

\hypertarget{learning-outcomes-52-learning-poverty-4-gepd-proficiency-in-grade-4}{%
\subsubsection{\texorpdfstring{\textbf{LEARNING OUTCOMES: 52\% LEARNING
POVERTY, 4\% GEPD PROFICIENCY IN GRADE
4}}{LEARNING OUTCOMES: 52\% LEARNING POVERTY, 4\% GEPD PROFICIENCY IN GRADE 4}}\label{learning-outcomes-52-learning-poverty-4-gepd-proficiency-in-grade-4}}

Learning poverty is defined as the share of children at end of primary
age below minimum reading proficiency, adjusted for out of school
children. Learning poverty in Jordan is 11\% points better than the
average for Middle East \& North Africa (excluding high income) region
and 23\% points worse than the average for Upper middle income
countries.

\hypertarget{figure-2.-leaning-poverty-comparison}{%
\paragraph{Figure 2. Leaning poverty
comparison}\label{figure-2.-leaning-poverty-comparison}}

\includegraphics[width=1\linewidth]{C:/Users/wb469649/Documents/Github/GEPD/Country_Reports/LP_LAYS_figures/lp_figure_JOR}

{\scriptsize
    \textcolor{darkgray}{\textit{Notes:} Grey circles represent other countries. Yellow circle represents Jordan. Other circles represent average learning poverty in Jordan's region and income group.}
  }

\textbf{GEPD grade 4 proficiency of 4\%} means 4\% students score
greater than 80\% in GEPD assessment. Student proficiency is 20 points
higher in language compared to numeracy, 2.1 points lower for boys
compared to girls, and 2 points lower in urban areas compared to rural
areas.

\hypertarget{figure-3.-gepd-grade-4-proficiency-jordan}{%
\paragraph{Figure 3. GEPD Grade 4 proficiency,
Jordan}\label{figure-3.-gepd-grade-4-proficiency-jordan}}

\includegraphics{C:/Users/wb469649/Documents/Github/GEPD/Country_Reports/Output/Jordan_2019_2Pager_files/figure-latex/learning_figure-1.pdf}

\hypertarget{comparing-de-facto-practices-and-policy-levers}{%
\subsubsection{\texorpdfstring{\textbf{COMPARING DE-FACTO PRACTICES AND
POLICY
LEVERS}}{COMPARING DE-FACTO PRACTICES AND POLICY LEVERS}}\label{comparing-de-facto-practices-and-policy-levers}}

Practice indicators measure quality of service delivery in schools such
as student performance, teacher knowledge, principal management skills,
etc. Policy lever indicators measure how well school, personnel and
student policies governing these practices are implemented. Comparing
de-facto practice and policy lever indicators allows identification of
low-scoring policy levers that affect observed practice indicators.

\hypertarget{teacher-effectiveness}{%
\paragraph{\texorpdfstring{\textbf{Teacher
effectiveness}}{Teacher effectiveness}}\label{teacher-effectiveness}}

Teacher content knowledge (34\%) needs improvement. Teacher proficiency
in language (25\%) is 33 points lower than mathematics proficiency
(58\%). Teacher pedagogical skills score (68\%) needs improvement, and
teacher attendance (93\%) is on target. Teaching - Support is the
weakest policy lever(2.6/5).

\begin{table}[H]
\resizebox{\columnwidth}{!}{\begin{tabular}{m{3.8cm}cm{4.2cm}c}
\multicolumn{2}{c}{\textbf{Practice Indicators}} & \multicolumn{2}{c}{\textbf{Policy levers (Teaching)}}\\\hline
Content knowledge                   & {\cellcolor{red!15}34\%}   & Attraction                                & \cellcolor{yellow!15}3.5 \\\cdashline{1-4} 
\hspace{1mm}\emph{Maths proficiency}               & {\cellcolor{red!15}58\%} &                                           & \cellcolor{yellow!15}\\\cdashline{1-2}   
\hspace{1mm}\emph{Language proficiency}            & {\cellcolor{red!15}25\%} & \multirow{-2}{*}{Selection \& deployment} & \multirow{-2}{*}{\cellcolor{yellow!15}3.5} \\\cdashline{1-4}        
Pedagogical skills                               & {\cellcolor{red!15}68\%}   &                                           & \cellcolor{red!15}\\\cdashline{1-2}
\hspace{1mm}\emph{\% Classroom culture}            & {\cellcolor{red!15}72\%} & \multirow{-2}{*}{Support}                 & \multirow{-2}{*}{\cellcolor{red!15}2.6} \\\cdashline{1-4}
\hspace{1mm}\emph{\% Instruction practices}        & {\cellcolor{red!15}74\%} &                                           & \cellcolor{yellow!15}\\\cdashline{1-2}
\hspace{1mm}\emph{\% Socio-emotional skills}       & {\cellcolor{red!15}68\%} & \multirow{-2}{*}{Evaluation}              & \multirow{-2}{*}{\cellcolor{yellow!15}3.6} \\\cdashline{1-4}
& \cellcolor{green!15} & Monitoring \& Accountability  & \cellcolor{yellow!15}3.1 \\\cdashline{3-4}
\multirow{-2}{*}{Teacher Attendance}             & \multirow{-2}{*}{\cellcolor{green!15}93\%} & Intrinsic motivation    & \cellcolor{yellow!15}3.4 \\\hline
\end{tabular}}
\\
\setstretch{0.8}\color{darkgray}\scriptsize{\textit{Notes:} Content knowledge(\& sub-indicators) indicate \% teachers with knowledge\textgreater{80\%}.Pedagogical skills(\& sub-indicators) indicate \% teachers with proficiency 3/5 or above.}
\end{table}

\hypertarget{capacity-for-learning-in-grade-1}{%
\paragraph{\texorpdfstring{\textbf{Capacity for learning in Grade
1}}{Capacity for learning in Grade 1}}\label{capacity-for-learning-in-grade-1}}

Student proficiency in Grade 1 (37\%) needs improvement. Socio-emotional
score(56)\% is the lowest knowledge sub-score. Student attendance(97\%)
is on target. Center-Based Care is the weakest policy lever(1.5/5).

\begin{table}[H]
\resizebox{\columnwidth}{!}{\begin{tabular}{m{3.8cm}cm{4.2cm}c}
\multicolumn{2}{c}{\textbf{Practice Indicators}} & \multicolumn{2}{c}{\textbf{Policy levers (Learners)}}\\\hline
Capacity for learning                  & {\cellcolor{red!15}37\%} & & \cellcolor{yellow!15}\\\cdashline{1-2}
\hspace{1mm}\emph{Numeracy score}        & \cellcolor{green!15}92   & \multirow{-2}{*}{Nutrition Programs} & \multirow{-2}{*}{\cellcolor{yellow!15}3.5}  \\\cdashline{1-4}  
\hspace{1mm}\emph{Literacy score}        & \cellcolor{red!15}76   & Health Programs                     & \cellcolor{yellow!15}3.6 \\\cdashline{1-4}  
\hspace{1mm}\emph{Executive score}       & \cellcolor{red!15}68   & Center based care                   & \cellcolor{red!15}1.5 \\\cdashline{1-4}  
\hspace{1mm}\emph{Socio-emotional score} & \cellcolor{red!15}56   & Caregiver Skills Capacity           & \cellcolor{green!15}4.7 \\\cdashline{1-4}  
Student Attendance                     & {\cellcolor{green!15}97\%}     & Caregiver Financial Capacity        & \cellcolor{green!15}4.4 \\\hline
\end{tabular}}
\\
\setstretch{0.8}\color{darkgray}\scriptsize{\textit{Notes:} Capacity for learning indicates \% students with knowledge\textgreater{80\%}. Sub-indicator scores refer to average subject knowledge on a 0-100 scale.}
\end{table}
\vfill\null

\hypertarget{inputs-infrastructure}{%
\paragraph{\texorpdfstring{\textbf{Inputs \&
Infrastructure}}{Inputs \& Infrastructure}}\label{inputs-infrastructure}}

Basic inputs (4.3/5) are on target. Percent of classrooms with a
functional blackboard and chalk(87)\% is the lowest score. Basic
infrastructure (4.1/5) is on target. Percent of schools accessible to
children with special needs(66)\% is the lowest score. Inputs \&
Infrastructure - Monitoring is the weakest policy lever(3/5).

\begin{table}[H]
\resizebox{\columnwidth}{!}{\begin{tabular}{m{3.8cm}cm{4.2cm}c}
\multicolumn{2}{c}{\textbf{Practice Indicators}} & \multicolumn{2}{c}{\textbf{Policy levers(Inputs)}}\\\hline
Basic inputs & \cellcolor{green!15}4.3 & & \cellcolor{yellow!15} \\\cdashline{1-2}
\hspace{1mm}\emph{\% Blackboard}    & {\cellcolor{yellow!15}87\%} & & \cellcolor{yellow!15}\\\cdashline{1-2}
\hspace{1mm}\emph{\% Stationery}    & {\cellcolor{green!15}98\%} & & \cellcolor{yellow!15}\\\cdashline{1-2}
\hspace{1mm}\emph{\% Furniture}     & {\cellcolor{yellow!15}88\%} & & \cellcolor{yellow!15}\\\cdashline{1-2}
\hspace{1mm}\emph{\% EdTech access} & {\cellcolor{yellow!15}89\%} & \multirow{-5}{4.2cm}{Inputs and infrastructure standards} & \multirow{-5}{*}{\cellcolor{yellow!15}3.8}\\\hline
Basic infrastructure & \cellcolor{green!15}4.1 & & \cellcolor{yellow!15} \\\cdashline{1-2}
\hspace{1mm}\emph{\% Drinking water}    & {\cellcolor{yellow!15}89\%} & & \cellcolor{yellow!15}\\\cdashline{1-2}
\hspace{1mm}\emph{\% Functional toilet} & {\cellcolor{red!15}69\%} & & \cellcolor{yellow!15}\\\cdashline{1-2}
\hspace{1mm}\emph{\% Internet}          & {\cellcolor{green!15}100\%} & & \cellcolor{yellow!15}\\\cdashline{1-2}
\hspace{1mm}\emph{\% Electricity}       & {\cellcolor{yellow!15}88\%} & & \cellcolor{yellow!15}\\\cdashline{1-2}
\hspace{1mm}\emph{\% Disability access} & {\cellcolor{red!15}66\%} & \multirow{-6}{4.2cm}{Inputs and infrastructure monitoring} & \multirow{-6}{*}{\cellcolor{yellow!15}3}\\\hline
\end{tabular}}
\\

    {\scriptsize
    \textcolor{darkgray}{\textit{Notes:} \% refers to \% schools with the given sub-component}
  }

\end{table}

\hypertarget{school-management-by-principals}{%
\paragraph{\texorpdfstring{\textbf{School Management by
principals}}{School Management by principals}}\label{school-management-by-principals}}

In school management, the lowest score is for principal's Instructional
Leadership(3.3/5), whereas the highest score is obtained for Operational
Management(4.7/5). School Management- Support is the weakest policy
lever(3.5/5).

\begin{table}[H]
\resizebox{\columnwidth}{!}{\begin{tabular}{m{4cm}cm{4cm}c}
\multicolumn{2}{c}{\textbf{Practice Indicators}} & \multicolumn{2}{c}{\textbf{Policy levers(Management)}}\\\hline
Operational management                       & \cellcolor{green!15}4.7                   &                                       & \cellcolor{green!15}\\\cdashline{1-2}
Instructional Leadership                     & \cellcolor{yellow!15}3.3                   & \multirow{-2}{*}{Clarity of functions} & \multirow{-2}{*}{\cellcolor{green!15}5} \\\hline
                                             & \cellcolor{yellow!15}                    & Attraction                            & \cellcolor{green!15}4.5\\\cdashline{3-4}
\multirow{-2}{*}{Principal School knowledge}  & \multirow{-2}{*}{\cellcolor{yellow!15}3.8} & Selection \& Deployment               & \cellcolor{green!15}4.6\\\cdashline{3-4}
                                              & \cellcolor{green!15}                    & Support                               & \cellcolor{yellow!15}3.5\\\cdashline{3-4}
\multirow{-2}{*}{Principal Management skills}  & \multirow{-2}{*}{\cellcolor{green!15}4} & Evaluation                            & \cellcolor{yellow!15}3.6\\\hline
\end{tabular}}
\\

    {\scriptsize
    \textcolor{darkgray}{\textit{Notes:} All scores are on a (0-5) scale.}
  }

\end{table}

\hypertarget{politics-bureaucratic-capacity-indicators}{%
\subsubsection{\texorpdfstring{\textbf{POLITICS \& BUREAUCRATIC CAPACITY
INDICATORS}}{POLITICS \& BUREAUCRATIC CAPACITY INDICATORS}}\label{politics-bureaucratic-capacity-indicators}}

Politics and bureaucratic capacity indicators measure the capacity and
orientation of the bureaucracy, as well as political factors affecting
education outcomes. The highest score in politics and bureaucratic
capacity is noted for Financing (4.7/5), and the lowest score is noted
for Impartial Decision-Making (3.1/5).

\begin{table}[H]
\resizebox{\columnwidth}{!}{\begin{tabular}{p{0.4\textwidth}c}
\textbf{Indicator} & \textbf{Value}\\ \hline
Quality of Bureaucracy & \cellcolor{yellow!15}3.6\\\cdashline{1-2}
Impartial Decision-Making & \cellcolor{yellow!15}3.1\\\cdashline{1-2}
Mandates \& Accountability & \cellcolor{green!15}4.2\\\cdashline{1-2}
National Learning Goals & \cellcolor{yellow!15}3.6\\\cdashline{1-2}
Financing & \cellcolor{green!15}4.7\\\hline
\end{tabular}}
\\

    {\scriptsize
    \textcolor{darkgray}{\textit{Notes:} All scores are on a (0-5) scale.}
  }

\end{table}
\raggedbottom

{\scriptsize
    \textcolor{darkgray}{\textit{Disclaimer:} GEPD numbers presented in this brief are based on multiple sources including GEPD instruments, UIS, GLAD and Learning Poverty indicators. For that reason, the numbers discussed here may be different from official statistics reported by governments and national offices of statistics. Such differences are due to the different purposes of the statistics, which can be for global comparison or to meet national definitions.}
  }

\end{document}
